\subsection{ANNIHILATORS. FAITHFUL MODULES. MONOGENOUS MODULES}

DEFINITION 11. The annihilator of a subset S of an A-module E is the set of elements \(\alpha \in \mathbf{A}\) such that \(\alpha x = 0\) for all \(x \in S\) .

The annihilator of S is usually denoted by Ann(S); for a subset S consisting of a single element x, we write Ann(x) instead of Ann(\{x\}) and call Ann(x) the annihilator of x.

The relation \(\alpha x = 0\) may also be expressed by saying that \(x\) is annihilated by \(\alpha\) .

It is immediate that the annihilator of an arbitrary subset S of E is a left ideal of A; for it to be equal to A, it is necessary and sufficient (by virtue of \((\mathbf{M}_{\mathrm{IV}})\) ) that \(S = \{0\}\) . If two subsets S, T of E are such that \(S \subset T\) , the annihilator of T is contained in the annihilator of S. If \((\mathbf{S}_{t})_{t \in I}\) is an arbitrary family of subsets of E, the annihilator of the union \(\bigcup_{t} S_{t}\) is the intersection of the annihilators of the \(S_{t}\) . In particular, the annihilator of a subset S of E is the intersection of the annihilators of the elements of S. To say that an element of E is free is equivalent to saying that its annihilator is \(\{0\}\) . For all \(x \in E\) and all \(\alpha \in A\) , the annihilator of \(\alpha x\) is the set of \(\beta \in A\) such that \(\beta \alpha \in \text{Ann}(x)\) .

The annihilator of a submodule M of E is a two-sided ideal of A; for, if \(\alpha x = 0\) for all \(x \in M\) , then also \(\alpha(\beta x) = 0\) for all \(x \in M\) and all \(\beta \in A\) , hence \(\alpha\beta\) belongs to the annihilator of M for all \(\beta \in A\) . In particular, the annihilator of E is a two-sided ideal of A.

For all \(\alpha \in \mathbf{A}\) , let \(h_{\alpha}\) be the homothety \(x \mapsto \alpha x\) ; it is known that the mapping \(\alpha \mapsto h_{\alpha}\) of \(\mathbf{A}\) into the endomorphism ring \(\mathcal{E} = \operatorname{Hom}_{\mathbf{Z}}(\mathbf{E}, \mathbf{E})\) of the commutative group (without operators) \(\mathbf{E}\) , is a ring homomorphism ( \(\S 2\) , no. 5). The inverse image of 0 under this homomorphism is the annihilator \(\mathfrak{a}\) of \(\mathbf{E}\) ; the image of \(\mathbf{A}\) under the homomorphism \(\alpha \mapsto h_{\alpha}\) is therefore isomorphic to the quotient ring \(\mathbf{A}/\mathfrak{a}\) . The module \(\mathbf{E}\) is called faithful if its annihilator \(\mathfrak{a}\) reduces to 0.

Let E be any A-module, a a two-sided ideal of A contained in Ann(E) and let \(\dot{\alpha}\) be an element of the quotient ring A/a; for all \(x \in E\) , the element \(\alpha x\) is the same for all the \(\alpha \in A\) belonging to the class \(\dot{\alpha} \mod a\) ; if this element is denoted by \(\dot{\alpha}x\) , it is immediately seen that the mapping \((\dot{\alpha}, x) \mapsto \dot{\alpha}x\) defines (with addition on E) an (A/a)-module structure on E. When \(a = \text{Ann}(E)\) , the (A/a)-module E thus defined is faithful; we shall say that it is the faithful module associated with the A-module E. Observe that every submodule of an A-module E is also a submodule of the associated faithful module and conversely.

DEFINITION 12. A module is called monogenous if it is generated by a single element.

Proposition 9 of no. 7 shows that, if E is a monogenous A-module and a is an element generating E, E is identical with the set A.a of \(\xi a\) , where \(\xi\) runs through A.

Examples. (1) Every monogenous group, being commutative (I, § 4, no. 10, Proposition 18), is a monogenous Z-module.

\begin{enumerate}
\def\labelenumi{(\arabic{enumi})}
\setcounter{enumi}{1}
\item
  If A is a commutative ring, the monogenous submodules of the A-module A are just the principal ideals (I, § 8, no. 6) of the ring A.
\item
  Every simple A-module E is monogenous, since the submodule of E generated by an element \(\neq0\) of E is necessarily equal to E.
\end{enumerate}

PROPOSITION 22. Let A be a ring. Every quotient module of \(\mathbf{A}_s\) is monogenous. Conversely, let E be a monogenous A-module, \(c\) a generator of E and a its annihilator; the linear mapping \(\xi \mapsto \xi c\) defines, when passing to the quotient, an isomorphism of \(\mathbf{A}_s / \mathfrak{a}\) onto E.

As \(A_{s}\) is itself monogenous, being generated by 1, the first assertion follows from no. 7, Corollary 1 to Proposition 9. The second is obvious, since \(\xi \mapsto \xi c\) is by hypothesis surjective and has kernel a.

Note that, if A is not commutative, the annihilators of two distinct generators c, \(c'\) of a monogenous A-module E are in general distinct and are also distinct from the annihilator of the module E. On the other hand, if A is commutative, the annihilator of a generator c of E is contained in the annihilator of every element of E and hence is the annihilator of the whole of E.

COROLLARY. Every submodule of a monogenous A-module E is isomorphic to a quotient module b/a where a and b are two left-ideals of A such that \(a \subset b\) . Every quotient module of a monogenous A-module is monogenous.

The second assertion is immediate and the first follows from Proposition 22 and I, § 4, no. 6, Theorem 4.

Note on the other hand that a submodule of a monogenous module is not necessarily monogenous. For example, if A is a commutative ring in which there exist non-principal ideals (VII, §1, no. 1), these ideals are non-monogenous submodules of the monogenous A-module A.

It follows from the definitions that a submodule of an A-module E generated by a family \((a_{t})\) of elements of E is the sum of the monogenous submodules

\(Aa_{t}\) of E; for \((a_{t})\) to be a basis of E, it is necessary and sufficient that each of the \(a_{t}\) be a free element of E and that the sum of the \(Aa_{t}\) be direct.

PROPOSITION 23. Let E be an A-module, the direct sum of an infinite family \((\mathbf{M}_t)_{t\in I}\) of non-zero submodules. For every generating system S of E, \(\operatorname{Card}(\mathbf{S})\geqslant \operatorname{Card}(\mathbf{I})\)

For all \(x \in S\) , let \(C_x\) be the finite set of indices \(i \in I\) such that the component of \(x\) in \(M_i\) is \(\neq 0\) and let \(C = \bigcup_{x \in S} C_x\) . Every \(x \in S\) belongs by definition to the submodule of \(E\) the direct sum of the \(M_i\) for \(i \in C\) and the hypothesis that \(S\) generates \(E\) implies therefore that \(C = I\) ; as \(I\) is by hypothesis infinite, so is \(S\) (Set Theory, III, §5, no. 1, Corollary 1 to Proposition 1); therefore \(\text{Card}(I) = \text{Card}(C) \leqslant \text{Card}(S)\) (Set Theory, III, §6, no. 3, Corollary 3 to Theorem 2).

COROLLARY 1. Under the hypotheses of Proposition 23, suppose that each \(\mathbf{M}_{\iota}\) is monogenous and that \(\mathbf{E}\) is the direct sum of a second family \((\mathrm{N}_{\lambda})_{\lambda \in \mathbf{L}}\) of non-zero monogenous submodules. Then \(\operatorname{Card}(\mathbf{L}) = \operatorname{Card}(\mathbf{I})\) .

If \(b_{\lambda}\) is a generator of \(\mathbf{N}_{\lambda}\) , the set of \(b_{\lambda}\) is a generating system of \(\mathbf{E}\) , hence \(\operatorname{Card}(\mathbf{L}) \geqslant \operatorname{Card}(\mathbf{I})\) . In particular \(\mathbf{L}\) is infinite and, exchanging the roles of \((\mathbf{M}_i)\) and \((\mathbf{N}_{\lambda})\) , similarly \(\operatorname{Card}(\mathbf{I}) \geqslant \operatorname{Card}(\mathbf{L})\) , whence the corollary.

COROLLARY 2. If a module E admits an infinite basis B, every generating system of E has cardinal \(\geqslant\) Card(B) and every basis of E is equipotent to B.

